\section{练习}

\begin{question}\begin{drill}[落球在运动车面上的摩擦冲量]\label{em5:drill:1}
将小球视为不计自转的质点。质量$M$的小车以速率$v$向左运动，质量$m$的小球由高$h$处自由落下，碰车后反弹到相同高度。碰撞中支持力远大于重力，接触动摩擦因数$\mu$，若相对水平速度消失后不再滑动。求球的水平末速度。
\end{drill}
\fig{fig:em5:drill:1:s}{题设：落球在运动车面上的摩擦冲量}{\begin{tikzpicture}[scale=.9]\draw[world](0,0)rectangle(5,.5);\draw[flow](1,-.4)--(-.3,-.4)node[left]{$v$};\fill(3,2.5)circle(.1);\draw[flow](3,2.3)--(3,1.1)node[right]{$\sqrt{2gh}$};\draw[dim](4,.5)--(4,2.5)node[midway,right]{$h$};\end{tikzpicture}}
\end{question}
\answerlink{em5:drill:1}
\begin{question}\begin{drill}[桌边麦管与两只苍蝇]\label{em5:drill:2}
均匀细杆质量$M$、长$L$，初始一半伸出光滑水平桌边，另一半在桌内。质量$m<M$的苍蝇从桌内杆端缓慢爬到桌外端，杆始终水平且不倾覆。随后另一只质量$m^{\prime}$的苍蝇缓慢竖直降落到第一只身上，着陆冲击可忽略。忽略水平冲量，求仍不倾覆的$m^{\prime}$上限。
\end{drill}
\fig{fig:em5:drill:2:s}{题设：桌边麦管与两只苍蝇}{\begin{tikzpicture}[scale=.9]\draw[world](-3,0)--(0,0)--(0,-1.5);\draw[very thick](-2.5,.15)--(2.5,.15);\fill(-2.5,.4)circle(.08)node[above]{$m$};\draw[guide](0,-.2)--(0,1.8);\node at(1.8,-.4){桌外};\draw[flow](-2,1.2)--(1,1.2);\end{tikzpicture}}
\end{question}
\answerlink{em5:drill:2}
\begin{question}\begin{drill}[叠放弹性小球的反弹]\label{em5:drill:3}
两小球沿同一竖线下落，下面质量$M$、上面质量$m$（$M>m$）。从高$h$处静止释放，忽略两球尺寸，假定下球先与地面完全弹性碰撞，再立即与仍向下的上球完全弹性碰撞。求上球首次弹起最高高度，并取$M\gg m$极限。
\end{drill}
\fig{fig:em5:drill:3:s}{球间距作放大示意，计算忽略球尺寸及两次碰撞之间的高度变化}{\begin{tikzpicture}

\draw[world](-1,0)--(3,0);\fill(.5,2.5)circle(.18)node[right]{$M$};\fill(.5,2.88)circle(.1)node[right]{$m$};
\draw[flow](1.7,2.8)--(1.7,1.6)node[right]{$u$};\draw[dim](-.4,0)--(-.4,2.5)node[midway,left]{$h$};\node at(.8,3.5){由静止释放};
\end{tikzpicture}}
\end{question}
\answerlink{em5:drill:3}
